\chapter{New Iterator and Algorithm Facilities}
\label{ch:c14-iterator-algorithm}

\begin{quote}
\emph{C++11 made generic, range-based iteration the default style; C++14 patched the rough edges that style exposed. The non-member \texttt{c}/\texttt{r} range accessors complete the free-function family, \texttt{make\_reverse\_iterator} removes a type-spelling chore, value-initialized forward iterators gain defined comparison so they can serve as sentinels, and the two-range \texttt{<algorithm>} overloads close a long-standing buffer-overrun hole.}
\end{quote}

These additions share a theme: they make iterator-and-algorithm code that \emph{looked} correct in C++11 actually safe and uniform. The headline item is the new four-iterator overloads of \texttt{equal}, \texttt{mismatch}, and \texttt{is\_permutation}, which finally let the standard algorithms know where the \textbf{second} range ends --- eliminating a class of out-of-bounds reads that the three-iterator forms invite. The rest are ergonomic completions of the C++11 generic-iteration model.

This chapter is authored to complete the C++14 coverage; it has no predecessor source material.

\section{Non-Member \texttt{cbegin}/\texttt{cend}/\texttt{rbegin}/\texttt{rend}/\texttt{crbegin}/\texttt{crend}}
\label{sec:c14-nonmember-accessors}

C++11 added the non-member \texttt{std::begin(c)} and \texttt{std::end(c)} so generic code could iterate any range --- including C arrays --- through a uniform free-function call. But it stopped there: the \texttt{const} and reverse accessors existed only as \emph{member} functions, so generic code that wanted a const or reverse iterator had to assume a member interface, which arrays do not have. \textbf{C++14 completed the set} with non-member \texttt{std::cbegin}, \texttt{std::cend}, \texttt{std::rbegin}, \texttt{std::rend}, \texttt{std::crbegin}, and \texttt{std::crend}.

\begin{lstlisting}[language=C++,caption={Uniform const/reverse access across containers AND arrays}]
#include <iterator>
#include <vector>

template <typename Range>
auto sum_const(const Range& r) {
    typename Range::value_type acc{};
    for (auto it = std::cbegin(r); it != std::cend(r); ++it)
        acc += *it;                 // *it is read-only regardless of r's constness
    return acc;
}

std::vector<int> v{1, 2, 3};
int  arr[] = {4, 5, 6};

auto a = std::cbegin(v);            // const_iterator into the vector
auto b = std::cbegin(arr);         // const int* -- works on a raw array too
auto r = std::crbegin(v);          // const reverse iterator, last element first
\end{lstlisting}

The win is \textbf{uniformity}: a template can now obtain a const or reverse iterator from \emph{any} range --- STL container, \texttt{std::array}, or built-in array --- through one spelling, with no \texttt{typename Range::const\_iterator} and no member-function assumption. \texttt{std::cbegin} also guarantees a true \texttt{const\_iterator} even when called on a non-\texttt{const} container.

\section{\texttt{std::make\_reverse\_iterator}}
\label{sec:c14-make-reverse-iterator}

Constructing a \texttt{std::reverse\_iterator} by hand requires naming the underlying iterator type, which in generic code means an ugly \texttt{std::reverse\_iterator<decltype(it)>(it)}. \textbf{C++14's \texttt{std::make\_reverse\_iterator(it)}} is a factory that \emph{deduces} that type --- the same convenience \texttt{make\_pair}/\texttt{make\_tuple} provide.

\begin{lstlisting}[language=C++,caption={Building reverse iterators without spelling the type}]
#include <iterator>
#include <vector>

std::vector<int> v{1, 2, 3, 4, 5};

// Verbose, C++11:
auto r1 = std::reverse_iterator<std::vector<int>::iterator>(v.end());

// C++14: the type is deduced from the argument:
auto r2    = std::make_reverse_iterator(v.end());     // points at 5
auto r2end = std::make_reverse_iterator(v.begin());
// [r2, r2end) traverses v back-to-front: 5, 4, 3, 2, 1
\end{lstlisting}

It is most valuable inside templates, where the iterator type is a dependent name:

\begin{lstlisting}[language=C++,caption={A generic ``find from the back'' using the factory}]
template <typename It, typename T>
It find_last(It first, It last, const T& value) {
    auto rfirst = std::make_reverse_iterator(last);
    auto rlast  = std::make_reverse_iterator(first);
    auto found  = std::find(rfirst, rlast, value);
    return found == rlast ? last : std::prev(found.base());
}
\end{lstlisting}

The factory carries the standard reverse-iterator semantics: a reverse iterator built from position \texttt{p} dereferences to \texttt{*(p-1)}, so \texttt{make\_reverse\_iterator(container.end())} points at the last element and \texttt{make\_reverse\_iterator(container.begin())} is the reverse end.

\section{Null (Value-Initialized) Forward Iterators}
\label{sec:c14-null-forward-iterators}

C++14 tightened the iterator requirements so that \textbf{value-initialized forward iterators may be compared, and all value-initialized iterators of the same type compare equal} --- regardless of which container (if any) they came from. Before this, a default-constructed forward iterator was \emph{singular}: comparing it was undefined behavior, so it could not portably be used as a sentinel.

\begin{lstlisting}[language=C++,caption={A value-initialized iterator is a well-defined ``null'' sentinel}]
#include <vector>

std::vector<int>::iterator a{};   // value-initialized -> null forward iterator
std::vector<int>::iterator b{};   // also null

bool same = (a == b);             // C++14: well-defined, and true
\end{lstlisting}

The guarantee makes a default-constructed iterator usable as a portable \textbf{``end of an empty range'' / not-found marker} --- the iterator analogue of \texttt{nullptr}.

\begin{lstlisting}[language=C++,caption={Using a null iterator as an explicit ``unset'' state}]
template <typename It>
class OptionalCursor {
    It pos_{};                              // null forward iterator = "unset"
public:
    bool engaged() const { return pos_ != It{}; }   // compare against the null value
    void set(It p) { pos_ = p; }
    It   get() const { return pos_; }
};
\end{lstlisting}

Two value-initialized iterators of the same type form an \textbf{empty range} \texttt{[It\{\}, It\{\})}, which any algorithm processes as zero elements. The contract is narrow but important: it applies to \emph{value-initialized} iterators, not to arbitrary singular iterators left over from a destroyed container --- those remain invalid.

\section{Two-Range \texttt{<algorithm>} Overloads}
\label{sec:c14-two-range-algorithms}

This is the most consequential addition in the chapter. C++11's \texttt{std::equal}, \texttt{std::mismatch}, and \texttt{std::is\_permutation} came in a \textbf{three-iterator} form: \texttt{(first1, last1, first2)}. They receive the end of the \emph{first} range but only the \emph{beginning} of the second, and simply assume the second range is at least as long. If it is shorter, the algorithm reads past its end --- \textbf{undefined behavior}.

\begin{lstlisting}[language=C++,caption={The three-iterator form is unsafe when ranges differ in length}]
#include <algorithm>
#include <vector>

std::vector<int> a{1, 2, 3, 4, 5};
std::vector<int> b{1, 2, 3};

// THREE-iterator form: reads b[3], b[4] -- PAST THE END of b. UB.
bool bad = std::equal(a.begin(), a.end(), b.begin());   // do NOT do this
\end{lstlisting}

\textbf{C++14 added four-iterator overloads} --- \texttt{(first1, last1, first2, last2)} --- that take the end of the second range too. They compare both lengths and never read past either end: \texttt{equal} returns \texttt{false} immediately if the ranges differ in length, \texttt{mismatch} stops at whichever range ends first, and \texttt{is\_permutation} knows both sizes up front.

\begin{lstlisting}[language=C++,caption={The C++14 four-iterator form is length-safe}]
#include <algorithm>
#include <vector>

std::vector<int> a{1, 2, 3, 4, 5};
std::vector<int> b{1, 2, 3};

bool eq = std::equal(a.begin(), a.end(), b.begin(), b.end());  // false: sizes differ
                                                               // and NO overrun

auto m = std::mismatch(a.begin(), a.end(), b.begin(), b.end());
// stops at b.end(); m.first points into a at the first unmatched/leftover position

std::vector<int> p{3, 1, 2};
bool perm = std::is_permutation(b.begin(), b.end(), p.begin(), p.end());  // true
\end{lstlisting}

Each algorithm also has a four-iterator overload taking a custom binary predicate:

\begin{lstlisting}[language=C++,caption={Four-iterator form with a predicate}]
bool ci_equal = std::equal(s1.begin(), s1.end(), s2.begin(), s2.end(),
                           [](char x, char y) {
                               return std::tolower(x) == std::tolower(y);
                           });
\end{lstlisting}

The three-iterator overloads were not removed --- they remain for the case where you have \emph{already} verified the lengths match --- but the four-iterator forms should be your default. They turn a latent buffer overrun into a correct \texttt{false}.

\begin{quote}
\textbf{Godhood tip:} treat the three-iterator \texttt{equal}/\texttt{mismatch}/\texttt{is\_permutation} as a code smell in review. Unless a comment proves the second range's length is guaranteed, the four-iterator form is the only safe choice --- and it is also the one that gives the right answer for unequal lengths instead of UB.
\end{quote}

\section{Professional Insights}
\label{sec:c14-iterator-algorithm-insights}

\textbf{Use the non-member \texttt{c}/\texttt{r} accessors in every generic algorithm.} \texttt{std::cbegin(r)}/\texttt{std::cend(r)} give you a read-only traversal of \emph{any} range --- container or raw array --- without assuming a member interface or spelling a dependent \texttt{const\_iterator} type. Defaulting to them makes template code work uniformly across more argument types.

\textbf{Prefer \texttt{make\_reverse\_iterator} to a hand-spelled \texttt{reverse\_iterator<...>}.} In generic code the underlying iterator is a dependent type; the factory deduces it, keeping reverse-traversal helpers short and refactor-stable.

\textbf{A value-initialized iterator is the portable iterator equivalent of \texttt{nullptr}.} When you need an ``unset position'' or an empty-range sentinel in generic code, value-initialize the iterator type and compare against \texttt{It\{\}} --- C++14 makes that comparison defined. Don't, however, treat \emph{any} singular iterator as comparable; the guarantee is specifically for value-initialized ones.

\textbf{Make the four-iterator \texttt{equal}/\texttt{mismatch}/\texttt{is\_permutation} your default --- always.} The three-iterator forms silently read past the end of the second range when lengths differ, which is undefined behavior and a genuine security and stability bug. Passing both ends costs nothing, removes the overrun, and returns the correct result for unequal-length inputs. In latency-critical and security-sensitive code especially, the four-iterator overload is the only defensible choice.
